\chapter{Logical View \\
\small{\textit{-- Ryan Davis, Roy Tas, and Cory Vitanza}}
\index{Logical View} 
\index{Chapter!Logical View}
\label{Chapter::LogicalView}}

Class diagrams present the structure of our Design, highlighting the firmware on our Wio Tracker 1110, and how that data is processed in our React software.
%class diagram good

\begin{figure}[h]
\centering
\includegraphics[width=0.7\textwidth]{png/diagrams/Firmware Class Diagram.jpg}
\caption{\FigureLabel{FirmwareClasses}: Class diagram for the Wio Tracker 1110 Development Board's Firmware}
\label{fig:FirmwareClasses}
\end{figure}

\begin{figure}[h]
\centering
\includegraphics[width=1.15\textwidth]{png/diagrams/ReactMVC.jpg}
\caption{\FigureLabel{Model View Controller}: Class Diagram for our React Interface Using the model view controller design pattern.}
\label{fig:MVCClasses}
\end{figure}

\section*{Class Diagram Description}

\begin{itemize}

    \item \textbf{Overview}
    \begin{itemize}
        \item The class diagram is organized into four major packages: \textbf{Entities}, \textbf{Models}, \textbf{Controllers}, and \textbf{Views}.
        \item This structure follows the Model--View--Controller (MVC) architectural pattern.
        \item Each package contains classes with specific responsibilities and well-defined relationships such as inheritance, composition, realization, and dependencies.
    \end{itemize}

    \item \textbf{Entities Package}
    \begin{itemize}
        \item Contains raw data structures corresponding to the Supabase database tables.
        \item Includes \texttt{BoardEntity}, \texttt{GeofenceEntity}, \texttt{AccountEntity}, and \texttt{LocationRecordEntity}.
        \item All entities implement the \texttt{IEntity} interface, which standardizes the attribute \texttt{id} across all stored records.
        \item Entities contain no domain logic and only represent database rows.
    \end{itemize}

    \item \textbf{Models Package}
    \begin{itemize}
        \item Models wrap the Entities through composition and extend the abstract \texttt{Model} base class.
        \item Responsible for data cleaning, validation, and transformation (e.g., converting timestamps to \texttt{Date} objects).
        \item Contain computed attributes required by the UI, such as geofence containment, distance calculations, battery warnings, and formatted display information.
        \item Includes \texttt{BoardModel}, \texttt{GeofenceModel}, \texttt{AccountModel}, and \texttt{LocationRecordModel}.
    \end{itemize}

    \item \textbf{Controllers Package}
    \begin{itemize}
        \item All controllers implement the \texttt{Controller} interface.
        \item Responsible for orchestrating system logic, including loading database data, updating models, applying filters, and managing CRUD operations.
        \item Controllers detect events such as geofence breaches or location abnormalities, and trigger notifications via \texttt{NotificationService}.
        \item Includes \texttt{BoardController}, \texttt{GeofenceController}, \texttt{FilteringController}, and \texttt{NotificationController}.
    \end{itemize}

    \item \textbf{Views Package}
    \begin{itemize}
        \item Views depend on their corresponding controller to retrieve and update models.
        \item Views do not contain domain logic---their only responsibility is UI rendering.
        \item Views interact with the \texttt{MapComponent} to display board locations, geofences, filters, and violations.
        \item Includes \texttt{BoardView}, \texttt{GeofenceView}, \texttt{AccountView}, and the shared \texttt{MapComponent}.
    \end{itemize}

    \item \textbf{Key Relationships}
    \begin{itemize}
        \item \textbf{Inheritance:}
        \begin{itemize}
            \item Model subclasses extend the abstract \texttt{Model}.
            \item Controllers implement \texttt{IController}.
            \item Views depend on \texttt{MapView}.
        \end{itemize}
        \item \textbf{Composition:}
        \begin{itemize}
            \item Each Model class owns its corresponding Entity (e.g., \texttt{BoardModel} has a \texttt{BoardEntity}).
        \end{itemize}
        \item \textbf{Dependencies:}
        \begin{itemize}
            \item Views depend on Controllers.
            \item Controllers depend on Models and backend API calls.
            \item Views depend on the shared \texttt{MapComponent} for UI rendering.
        \end{itemize}
        \item \textbf{Separation of Concerns:}
        \begin{itemize}
            \item Views never access Entities or perform domain logic.
            \item Models never render UI.
            \item Controllers manage logic but never render UI directly.
        \end{itemize}
    \end{itemize}

    \item \textbf{Summary}
    \begin{itemize}
        \item The diagram demonstrates a complete and maintainable MVC architecture.
        \item Entities represent raw data, Models transform and prepare information for use, Controllers coordinate and implement CRUD, and Views render the user-facing interface.
        \item This architecture supports geofence and board management, filtering operations, map visualization, and notification handling in the FindIt system.
    \end{itemize}

\end{itemize}
